% ===================================================================
%  TABELO N° 13 — La Hotelo. --- La Restorerio. La Kafeerio.
%  Feuillets 87 a 90 du fac-simile = folios 85 a 88.
%  Correspondance : folio imprime = numero de feuillet - 2.
%
%  Ca tabelo aparteni a un narativo a la unesma persono (« me »), sen
%  scenizado dialogala. Ol portas DU noti d'autoro : la unesma, folio
%  85, expliktas ke la detali di « chambro » ja aparis en la tabelo 6 ;
%  la duesma, folio 87, atencionigas ke la menuo sequas la disho-listo
%  dil vortaro. Ambwa noti aparas en la fac-simile en gras-italika —
%  divergante de la stilo (romano simpla) di la noti dil tabelo 1 e 5 —
%  e la transskripto sequas fidele ico.
% ===================================================================

% --- Feuillet 87 (folio 85, apertajo di la tabelo) ------------------
\begin{VUpage}[87]{85}
\VUnotes{143.2mm}{%
\textit{\VUgras{(*) Pri la detali di chambro, videz la tabelo N\textsuperscript{o}. 6.}}%
}

%%K t13-apar-1 apar
\VUsaut{3.0mm}
\VUtitre{15.0pt}{60}{TABELO N\textsuperscript{o} 13}
\VUsaut{1.6mm}
\VUpk{10.2pt}{40}{La Hotelo. --- La Restorerio.}
\VUsaut{2.0mm}
\VUpk{10.2pt}{40}{La Kafeerio.}
\VUsaut{3.4mm}

%%K t13-01-1 p
\VUblancAlinea
1. --- Arivinte hiere vespere en Royan, me\nl
lojeskis che « \textit{Hotelo di Oceano} » qua esabis re\cc
komendata a me da amiko.

%%K t13-01-2 p
\VUblancAlinea
On furnisis a me bela \VUgras{chambro} \textsuperscript{(2)}, sur la\nl
duesma \VUgras{etajo} \textsuperscript{(1)} ube me tre bone dormis (*).

%%K t13-02-1 p
\VUblancAlinea
2. --- Camatine, ekiranta mea chambro, en\nl
qua la \VUgras{servistino} \textsuperscript{(3)} adportabis la dejuneto,\nl
me eniris la \VUgras{fumeyo} \textsuperscript{(5)} qua anke esas uzata\nl
kom \VUgras{lekteyo} \textsuperscript{(4)} por lektar la \VUgras{jurnali} \textsuperscript{(6)}, fu\cc
mante \VUgras{sigareto} \textsuperscript{(8)}. Pos lektir, me decensis per\nl
l'\VUgras{acensilo} \textsuperscript{(9)} ed iris promenar tra la urbo.

%%K t13-03-1 p
\VUblancAlinea
3. --- Pro ke me obliviabis questionar la\nl
\VUgras{kontoristo} \textsuperscript{(12)} di la hotelo pri qua kloko on\nl
donas la repasti sur la \VUgras{komuna tablo} \textsuperscript{(11)}, me\nl
frue retrovenis e profitis lo por irar balnar en\nl
la \VUgras{balneyo} \textsuperscript{(10)} di la hotelo. Pose me iris itere\nl
a la \VUgras{kaseyo} \textsuperscript{(15)} por telefonar ad un de mea\nl
amiki. Ma la \VUgras{telefonilo} \textsuperscript{(13)} esis okupata; vi\cc
dante mea embaraseso la \VUgras{kasistino} \textsuperscript{(14)}, qua\nl
esis registraganta \VUgras{fakturo} \textsuperscript{(17)} sur la \VUgras{voyajero}\cc
\VUgras{registro} \textsuperscript{(16)} pregis la \VUgras{pordistulo} \textsuperscript{(18)} sendar la\nl
\VUgras{grumo} \textsuperscript{(19)} facar mea komiso \VUgras{bicikle} \textsuperscript{(20)}.

%%K t13-04-1 p
\VUblancAlinea
4. --- Me dankis elu ed ekiris por donar gra\cc
tifiketo a la \VUgras{portisto} \textsuperscript{(24)} (la penlaboristo) qua
\parplein
\end{VUpage}

% --- Feuillet 88 (folio 86) ------------------------------------------
\begin{VUpage}[88]{86}
%%K t13-04-1 p suite
\VUcontinue
adportis de la staciono sur sua \VUgras{manuvetu}\cc
\VUgras{ro} \textsuperscript{(25)} mea \VUgras{chapel-buxo} \textsuperscript{(26)}, \VUgras{mea valizo} \textsuperscript{(27)} e\nl
\VUgras{mea voyaj-sako} \textsuperscript{(28)}. La \VUgras{letristo} \textsuperscript{(21)} jus ari\cc
vinta, donis a me \VUgras{letro} \textsuperscript{(22)}.

%%K t13-05-1 p
\VUblancAlinea
5. --- Me restis ankore kelkatempe sur la\nl
solio dil pordo, proxim la \VUgras{letrobuxo} \textsuperscript{(23)}, re\cc
gardonte la cirkulado sur la strado. La \VUgras{kon}\cc
\VUgras{duktoro} \textsuperscript{(30)} dil \VUgras{omnibuso} \textsuperscript{(29)} esis karganta sur\nl
sua veturo la \VUgras{kofro} \textsuperscript{(31)} di \VUgras{voyajero} \textsuperscript{(32)} quan\nl
la \VUgras{hotel-mastro} \textsuperscript{(33)} esis adianta. Yuna \VUgras{tele}\cc
\VUgras{grafisto} \textsuperscript{(35)} senhaste adportis \VUgras{telegramo} \textsuperscript{(34)}\nl
por kliento di la hotelo; \VUgras{turisto} \textsuperscript{(36)} kun sua\nl
\VUgras{fotograf-aparato} \textsuperscript{(38)} bandoliere, konsultis sua\nl
\VUgras{guid-libro} \textsuperscript{(37)}.

%%K t13-tit-1 sub
\VUsaut{4.0mm}
\VUpk{10.2pt}{40}{Dejuno-tempo.}
\VUsaut{3.0mm}

%%K t13-06-1 p
\VUblancAlinea
6. --- Avan la greto, sensucia \VUgras{komisionis}\cc
\VUgras{to} \textsuperscript{(39)} dormetis inter sua \VUgras{porto-hoko} \textsuperscript{(40)} e sua\nl
\VUgras{ciraj-kesto} \textsuperscript{(41)}, dum ke yuna \VUgras{(linjo-)lavisti}\cc
\VUgras{no} \textsuperscript{(42)} qua portis \VUgras{korbo} \textsuperscript{(43)} de linjaro, ridis pro\nl
la solena mieno dil \VUgras{tabl-oficisto} \textsuperscript{(44)} qua stacis\nl
apud la pordo.

%%K t13-07-1 p
\VUblancAlinea
7. --- Vidante la \VUgras{koquistulo} \textsuperscript{(45)}, qua ekiris\nl
sua \VUgras{koqueyo} \textsuperscript{(47)} situita ye la \VUgras{sub-ter-etajo} \textsuperscript{(48)}\nl
por sendar \VUgras{skuristo} \textsuperscript{(46)} facar komiso, me me\cc
moris ke la horo dil dejuno esas balda e me\nl
eniris la restoreyo, trairante la korto en qua\nl
\VUgras{motoristo} \textsuperscript{(50)} garis sua \VUgras{automobilo} \textsuperscript{(49)}, dum ke\nl
\VUgras{mashinisto} \textsuperscript{(52)} reparis, segun l'indiki dil \VUgras{ci}\cc
\VUgras{klisto} \textsuperscript{(53)}, \VUgras{petrol-triciklo} \textsuperscript{(51)} qua jus esis aci\cc
dentita.

%%K t13-08-1 p
\VUblancAlinea
8. --- Me questionis yuna \VUgras{grumo} \textsuperscript{(54)} qua por\cc
tis \VUgras{bir-glasi} \textsuperscript{(55)}, ka la komuna tablo esas sub
\parplein
\end{VUpage}

% --- Feuillet 89 (folio 87) ------------------------------------------
\begin{VUpage}[89]{87}
\VUnotes{143.2mm}{%
\textit{\VUgras{(*) Per helpo di la disho-listo enduktita en la vortaro, la\nl
docisto darfos facile diversigar la suba menuo.}}%
}

%%K t13-08-1 p suite
\VUcontinue
la verando, e pro lua afirmanta respondo, me\nl
okupis plaso ye un de la tabli ed igis donar a\nl
me bonega repasto.

%%K t13-tit-2 sub
\VUsaut{4.0mm}
\VUpk{10.2pt}{40}{Bonega dejuno.}
\VUsaut{3.0mm}

%%K t13-09-1 p
\VUblancAlinea
9. --- La garsono adportis a me la menuo (*)\nl
dil dejuno e la vin-listo. Me selektis e komen\cc
dis kom \VUgras{acesorajo kreveti} kun \VUgras{butro,} e, kom\nl
entreo, \VUgras{tripi segun la modo di Caen.} Me ne\nl
manjis \VUgras{fisho,} ma demandis porciono de muto\cc
nyunal \VUgras{kruro,} e, kom \VUgras{legum-disho, spinati} kun\nl
kremo. Pose, pro ke nula de la \VUgras{entremesi} plezis\nl
a me, me igis adportar diversa deseri, \VUgras{fromajo,}\nl
\VUgras{frukti} e \VUgras{bisquiti.} Me pluse adjuntis bonega\nl
Saint-Emilion-\VUgras{vino,} e tre kontenta pri mea\nl
repasto me livis la tablo pos donacir al \VUgras{gar}\cc
\VUgras{sono} \textsuperscript{(57)} bel \VUgras{gratifiketo} \textsuperscript{(59)}.

%%K t13-10-1 p
\VUblancAlinea
10. --- Vidante me pronta por departar, la\nl
\VUgras{jerero} \textsuperscript{(60)} propozis a me adiror la balkono por\nl
admirar la splendida panoramo di \VUgras{Oceano} \textsuperscript{(62)}.\nl
Fore on videskis pesko-\VUgras{barketi} \textsuperscript{(63)}, \VUgras{seglo-na}\cc
\VUgras{vi} \textsuperscript{(64)} ed enorma \VUgras{paketboto} \textsuperscript{(65)} di qua la nigra\nl
silueto salias sur la blanka \VUgras{nubi} \textsuperscript{(67)} di la \VUgras{hori}\cc
\VUgras{zonto} \textsuperscript{(66)}. Lejera brizo igis fluktuar la \VUgras{stan}\cc
\VUgras{dardo} \textsuperscript{(70)} stekita super l'\VUgras{insigno} \textsuperscript{(69)} di la\nl
\VUgras{halo} \textsuperscript{(68)}.

%%K t13-tit-3 sub
\VUsaut{3.0mm}
\VUpk{10.2pt}{40}{Amuzi.}
\VUsaut{3.0mm}

%%K t13-11-1 p
\VUblancAlinea
11. --- La spektajo esis tante interesanta, ke\nl
me decidis acensor morge sur la teraso dil ka\cc
feerio kunportante \VUgras{bilorneto} \textsuperscript{(71)} o \VUgras{lorno} \textsuperscript{(72)}.
\end{VUpage}

% --- Feuillet 90 (folio 88) ------------------------------------------
\begin{VUpage}[90]{88}
%%K t13-12-1 p
\VUblancAlinea
12. --- Livante la balkono, me adiris la ka\cc
feerio di la hotelo por glutar \VUgras{drinkajo} \textsuperscript{(73)}, lu\cc
dante \VUgras{biliardo-partio} \textsuperscript{(75)}. Senhalte me trairis\nl
la kafe-salono en qua multa konsumanti ludis\nl
per \VUgras{shako} \textsuperscript{(78)}, per \VUgras{dami} \textsuperscript{(79)}, per \VUgras{domino-sto}\cc
\VUgras{ni} \textsuperscript{(80)}, per \VUgras{triktrako} \textsuperscript{(81)} o per \VUgras{karti} \textsuperscript{(82)}, e me\nl
direte acensis la \VUgras{biliardeyo} \textsuperscript{(74)}.

%%K t13-13-1 p
\VUblancAlinea
13. --- Kande la partio esis finita, me decen\cc
sis a la teretajo e demandis de la garsono \VUgras{ku}\cc
\VUgras{verto} \textsuperscript{(84)}, \VUgras{papero} \textsuperscript{(83)}, \VUgras{postal-marko} \textsuperscript{(85)}, \VUgras{plu}\cc
\VUgras{mo} \textsuperscript{(87)} ed \VUgras{inko} \textsuperscript{(86)} por skribar letro.

%%K t13-13-2 p
\VUblancAlinea
Pose me advokis \VUgras{veturisto} \textsuperscript{(92)} di qua la \VUgras{ve}\cc
\VUgras{turo} \textsuperscript{(88)} esis haltinta avan la kafeerio. Me ques\cc
tionis lu pri la preco segun la hori od un kuro,\nl
e, kande ni konkordis pri la preco, il apertis\nl
a me la \VUgras{vetur-pordo} \textsuperscript{(89)} ed instalis me. Il ria\cc
censis sur la \VUgras{sidilo} \textsuperscript{(91)}, rapide startigis la ka\cc
vali per \VUgras{flogeto} \textsuperscript{(93)} e ni departis por promeno\nl
charmanta.

\VUsaut{6.0mm}
\VUfilet{22mm}
\end{VUpage}
